% ASSIGNING WHAT IS ALREADY THERE IS A NO-OP -- IN eTeX MODE ONLY.
%
% This is the trap the method note warns about, caught in the act. Under
% the PRELOADED pdftex format, which is built with eTeX enabled, a group
% that assigns a value equal to the one standing gives nothing back:
%
%   \count5=7 {\count5=7}      no restore line
%   \count5=7 {\count5=8}      {restoring \count5=7}
%
% Under `pdftex -ini', which is what trip's oracle runs, BOTH draw a
% restore line. The same probe, the same engine, two answers -- the
% shortcut is an eTeX extension and is not in TeX.
%
%   pdftex -ini si.tex                  pdftex '\input si \end'
%   [A same value]                      [A same value]
%   {restoring \count5=7}               [B different value]
%   {restoring \dimen5=1.0pt}           {restoring \count5=7}
%   {restoring \catcode113=12}          [done]
%   [B different value]
%   {restoring \count5=7}
%   [done]
%
% HSTEX DOES THE TeX THING, because trip is what it is measured against
% and hstex has no eTeX mode to switch on. It was written the other way
% first, on a probe run under the preloaded format, and trip caught it:
% the reference gives back \fam and \displayindent from a display whose
% values were already what the display wanted, and the eTeX rule says it
% should not.
%
% If an eTeX mode is ever added, this is one of the things it turns on,
% along with `{reassigning ...}' and `{changing ...}' lines under
% \tracingassigns. What always saves either way is anything whose value
% is a freshly made thing and so never the same one twice: a glue, a
% token list, a \parshape, a box, a macro.
%
% RUN THIS ONE BOTH WAYS. As written it needs `pdftex -ini', which gives
% no catcodes of its own.
\catcode`\{=1 \catcode`\}=2 \catcode`\#=6
\tracingonline=1 \tracingrestores=1
\count5=7 \dimen5=1pt \catcode`\q=12
\message{[A same value]}
{\count5=7}
{\dimen5=1pt}
{\catcode`\q=12}
\message{[B different value]}
{\count5=8}
\message{[done]}
\end
